% ── Section 1: Introduction ───────────────────────────────────────────────────
% Responsible: B (xRFM person)
% Placeholder — to be filled by B

\section{Introduction}

% TODO (B): Write introduction covering:
%   - Challenges of tabular data prediction
%   - Motivation: why xRFM over GBDTs?
%   - Core idea of xRFM (AGOP + tree structure)
%   - Overview of experiments (A provides 1-2 sentences on datasets)
%   - Structure of the report

\textit{[Placeholder: Introduction to be written by B. Covers tabular data challenges, xRFM motivation, experimental overview.]}

Tabular data constitutes one of the most pervasive forms of data in industry and science \citep{grinsztajn2022tree}.
Despite dramatic advances in deep learning for images and text, Gradient Boosted Decision Trees (GBDTs) such
as XGBoost \citep{chen2016xgboost} and LightGBM \citep{ke2017lightgbm} remain the dominant approach for tabular
prediction tasks. Recently, \citet{beaglehole2025xrfm} proposed \textit{xRFM}, a model that combines
feature-learning kernel machines with an adaptive binary tree structure, achieving state-of-the-art results
on the TALENT benchmark.

In this report, we replicate the core experiments of \citet{beaglehole2025xrfm}, evaluating xRFM
against three baseline models — XGBoost, LightGBM, and MLP — across five tabular datasets spanning
regression and binary classification tasks. We further analyse scalability and interpretability to
contextualise the strengths and limitations of xRFM in practice.
